\section{Resultados: proteção}
\label{sec:protecao}

\subsection{Ataques impedidos por configuração (QP2)}

\PorFazer{Tabela $\pi_a(c)$: configurações em linhas, ataques em colunas, a partir
de \texttt{gerado/pi5/protection.csv}, com $S(c)$ na última coluna. Indicar o
número de repetições e os casos em que as repetições não concordaram.}

\PorFazer{Leitura: o que K1 já impede (tudo o que exige chegar ao \emph{broker});
o que só o envelope impede (K10 sem TLS, \emph{broker} ativo); o que só a
combinação impede (mensagem antiga nunca recebida: validade; repetição: janela);
o que só K8 e seguintes impedem (inundação, ocupação de ligações).}

\subsection{Cada ataque contra a sua defesa, com controlo negativo}

\PorFazer{Tabela dos cenários (\texttt{results/<campanha>/scenarios.md}): ataque,
defesa, previsão, resultado e evidência, com o controlo negativo ao lado. Esta
tabela confirma ou refuta, linha a linha, a tabela de ataques da Parte~1.}

\subsection{Observações}
\label{sec:observacoes}

Os ensaios mostraram efeitos que a tabela de ataques da Parte~1 não previa.
\PorFazer{Confirmar cada um com os dados do Raspberry~Pi e quantificar.}

\begin{itemize}[leftmargin=*]
  \item \textbf{Ocupação de ligações.} O limite de ligações protege os recursos do
    \emph{broker} e mantém os clientes já ligados, mas o atacante ocupa as vagas
    livres; se o \emph{broker} reiniciar durante o ataque, o atacante pode ficar
    com todas as vagas e nenhum sensor volta a ligar-se.
  \item \textbf{Chave roubada e janela da vítima.} As mensagens aceites com uma
    chave roubada avançam a janela anti-repetição dessa chave; o tráfego legítimo
    do dispositivo passa a ser rejeitado até à rotação de chaves.
  \item \textbf{Credencial clonada e sessão MQTT.} Se o atacante se liga com a
    identidade de um dispositivo que continua ligado, o \emph{broker} entrega-lhe
    a sessão e desliga o legítimo; os dois alternam enquanto o ataque dura, o que
    é uma negação de serviço da vítima.
  \item \textbf{Falsificação pelo recetor num \emph{broker} legítimo.} A ACL já
    impede um atuador de publicar nos tópicos de ordens dos outros; a diferença
    entre a chave partilhada e a assinatura só se observa sem ACL ou com um
    \emph{broker} ativo.
  \item \textbf{Atraso e supressão.} Um \emph{broker} ativo que atrase mensagens
    menos do que $\Delta_a$, ou que as suprima, não é detetado pelo recetor.
\end{itemize}
